\documentclass[11pt]{article}

\usepackage{amsmath,amssymb,amsthm}
\usepackage{geometry}
\usepackage{bm}
\usepackage[framemethod=TikZ]{mdframed}

% Table of Contents depth
\setcounter{tocdepth}{2}      % sections and subsections
\setcounter{secnumdepth}{3}   % number down to subsubsections

% Hyperlinks (load last)
\usepackage[
colorlinks=true,
linkcolor=blue,
citecolor=blue,
urlcolor=blue
]{hyperref}
\usepackage{bookmark}

\geometry{margin=1in}

\title{A Geometric Derivation of the Student's $t$ Distribution}
\author{}
\date{}

\begin{document}
	
	\maketitle
	\tableofcontents
	
	
	\section{Introduction}
	
	The classical definition of the Student's $t$ distribution is
	
	\[
	T=\frac{Z}{\sqrt{W/\nu}},
	\]
	
	where
	
	\[
	Z\sim N(0,1), \qquad
	W\sim\chi^2_\nu,
	\]
	
	and $Z$ and $W$ are independent.
	
	Although this construction is useful in statistics, it does not explain the
	geometric meaning of the distribution.
	
	The purpose of this derivation is to show that the Student's $t$
	distribution arises naturally from the geometry of high-dimensional
	spheres.
	
	The key result will be
	
	\[
	T
	=
	\sqrt{\nu}\cot\theta,
	\]
	
	where $\theta$ is the angle between a uniformly random point on the unit
	sphere $S^\nu$ and any fixed reference direction.
	
	%%%%%%%%%%%%%%%%%%%%%%%%%%%%%%%%%%%%%%%%%%%%%%%%%%%%%%%%%%%%%%%
	
	\section{Step 1. Constructing the Gaussian Vector}
	
	Let
	
	\[
	Z_1,Z_2,\ldots,Z_{\nu+1}
	\]
	
	be independent standard normal random variables.
	
	Define the vector
	
	\[
	\mathbf Z
	=
	(Z_1,\ldots,Z_{\nu+1})^T.
	\]
	
	Its joint density is
	
	\[
	f(\mathbf z)
	=
	(2\pi)^{-(\nu+1)/2}
	\exp
	\left(
	-\frac12
	\sum_{i=1}^{\nu+1}
	z_i^2
	\right).
	\]
	
	\[
	f(\mathbf z)
	=
	(2\pi)^{-(\nu+1)/2}
	e^{-\|\mathbf z\|^2/2}.
	\]
	
	\begin{mdframed}[
		linewidth=0.8pt,
		roundcorner=6pt,
		backgroundcolor=gray!5,
		linecolor=gray!60,
		innertopmargin=10pt,
		innerbottommargin=10pt,
		innerleftmargin=10pt,
		innerrightmargin=10pt
		]
		
		\subsection*{Derivation of the Joint Density of a Gaussian Vector}
		
		Let
		
		\[
		Z_1,Z_2,\ldots,Z_{\nu+1}
		\]
		
		be independent random variables, each following the standard normal
		distribution,
		
		\[
		Z_i\sim N(0,1),
		\qquad
		i=1,\ldots,\nu+1.
		\]
		
		Recall that the probability density function of a standard normal random
		variable is
		
		\[
		f_{Z_i}(z_i)
		=
		\frac{1}{\sqrt{2\pi}}
		\exp\left(-\frac{z_i^2}{2}\right).
		\]
		
		Since the random variables are assumed to be independent, the joint
		density is simply the product of the individual densities:
		
		\[
		f_{Z_1,\ldots,Z_{\nu+1}}
		(z_1,\ldots,z_{\nu+1})
		=
		\prod_{i=1}^{\nu+1}
		f_{Z_i}(z_i).
		\]
		
		Substituting the expression for each marginal density gives
		
		\[
		=
		\prod_{i=1}^{\nu+1}
		\left(
		\frac{1}{\sqrt{2\pi}}
		\exp\left(-\frac{z_i^2}{2}\right)
		\right).
		\]
		
		We now separate the product into two parts.
		
		The constant factors satisfy
		
		\[
		\prod_{i=1}^{\nu+1}
		\frac{1}{\sqrt{2\pi}}
		=
		\left(
		\frac{1}{\sqrt{2\pi}}
		\right)^{\nu+1}
		=
		(2\pi)^{-(\nu+1)/2}.
		\]
		
		Next, use the elementary identity
		
		\[
		\prod_{i=1}^{n}e^{a_i}
		=
		e^{\sum_{i=1}^{n}a_i}.
		\]
		
		Applying this to the exponential terms,
		
		\[
		\prod_{i=1}^{\nu+1}
		\exp\left(-\frac{z_i^2}{2}\right)
		=
		\exp\left(
		-\frac12
		\sum_{i=1}^{\nu+1}
		z_i^2
		\right).
		\]
		
		Combining both parts yields
		
		\[
		f_{Z_1,\ldots,Z_{\nu+1}}
		(z_1,\ldots,z_{\nu+1})
		=
		(2\pi)^{-(\nu+1)/2}
		\exp\left(
		-\frac12
		\sum_{i=1}^{\nu+1}
		z_i^2
		\right).
		\]
		
		Finally, define the random vector
		
		\[
		\mathbf Z
		=
		(Z_1,\ldots,Z_{\nu+1})^T,
		\]
		
		and let
		
		\[
		\mathbf z
		=
		(z_1,\ldots,z_{\nu+1})^T
		\]
		
		denote a realization of this vector.
		
		Using the Euclidean norm,
		
		\[
		\|\mathbf z\|^2
		=
		z_1^2+z_2^2+\cdots+z_{\nu+1}^2
		=
		\sum_{i=1}^{\nu+1}z_i^2,
		\]
		
		the joint density can be written compactly as
		
		\[
		\boxed{
			f_{\mathbf Z}(\mathbf z)
			=
			(2\pi)^{-(\nu+1)/2}
			\exp\left(
			-\frac12
			\|\mathbf z\|^2
			\right).
		}
		\]
		
		This expression reveals an important geometric property: the density
		depends on the vector $\mathbf z$ only through its Euclidean length
		$\|\mathbf z\|$. Consequently, every point lying on the same sphere
		centered at the origin has the same probability density.
		
		In other words, the level sets of the density are concentric spheres,
		
		\[
		\{\mathbf z:\|\mathbf z\|=r\},
		\]
		
		which immediately implies that the multivariate Gaussian distribution is
		\emph{rotationally invariant}. Rotating the vector about the origin does
		not change its density because orthogonal transformations preserve the
		Euclidean norm:
		
		\[
		\|Q\mathbf z\|
		=
		\|\mathbf z\|,
		\qquad
		Q^TQ=I.
		\]
		
		Therefore,
		
		\[
		f_{\mathbf Z}(Q\mathbf z)
		=
		f_{\mathbf Z}(\mathbf z),
		\]
		
		for every orthogonal matrix \(Q\), proving that the Gaussian vector is
		rotationally invariant.
		
	\end{mdframed}
	
	
	\section{Step 2. Polar Representation}
	
	Every rotationally invariant vector admits the decomposition
	
	\[
	\mathbf Z
	=
	R\mathbf U,
	\qquad 
	\text{where}
	\qquad 
	R=\|\mathbf Z\|,
	\] 
	
	\[ 
	\text{and}
	\qquad 
	\mathbf U
	=
	(U_1,\ldots,U_{\nu+1})
	\sim
	\mathcal{U}_{[S^\nu]}
	\]
	
	
	\[
	\text{where}
	\qquad S^\nu
	=
	\left\{
	x\in\mathbb R^{\nu+1}
	:
	\|x\|=1
	\right\}.
	\]
	
	\begin{mdframed}[
		linewidth=0.8pt,
		roundcorner=6pt,
		backgroundcolor=gray!5,
		linecolor=gray!60,
		innertopmargin=10pt,
		innerbottommargin=10pt,
		innerleftmargin=10pt,
		innerrightmargin=10pt
		]
		
		\subsection*{Polar Representation of a Rotationally Invariant Random Vector}
		
		The previous section showed that the Gaussian vector
		
		\[
		\mathbf Z=(Z_1,\ldots,Z_{\nu+1})^T
		\]
		
		has a density that depends only on its Euclidean norm,
		
		\[
		f_{\mathbf Z}(\mathbf z)
		=
		\phi(\|\mathbf z\|),
		\]
		
		for some function $\phi$.
		
		This means that every point lying on the same sphere centered at the origin
		has the same probability density.
		
		The natural question is therefore:
		
		\begin{center}
			\emph{Can every realization of $\mathbf Z$ be described by its distance
				from the origin and its direction?}
		\end{center}
		
		The answer is yes.
		
		\paragraph{Step 1. Define the Radius}
		
		Given any realization
		
		\[
		\mathbf z
		=
		(z_1,\ldots,z_{\nu+1})^T,
		\]
		
		its Euclidean length is
		
		\[
		R
		=
		\|\mathbf z\|
		=
		\sqrt{
			z_1^2+z_2^2+\cdots+z_{\nu+1}^2
		}.
		\]
		
		Since $R$ is a distance, we always have
		
		\[
		R\ge0.
		\]
		
		Thus every point in $\mathbb R^{\nu+1}$ is assigned a unique nonnegative
		radius.
		
		\vspace{0.3cm}
		
		\paragraph{Step 2. Define the Direction}
		
		Whenever $R>0$, divide the vector by its length:
		
		\[
		\mathbf U
		=
		\frac{\mathbf z}{R}.
		\]
		
		Writing this coordinatewise,
		
		\[
		\mathbf U
		=
		\left(
		\frac{z_1}{R},
		\frac{z_2}{R},
		\ldots,
		\frac{z_{\nu+1}}{R}
		\right)^T.
		\]
		
		We now verify that $\mathbf U$ is a unit vector.
		
		Indeed,
		
		\[
		\|\mathbf U\|^2
		=
		\left\|
		\frac{\mathbf z}{R}
		\right\|^2
		=
		\frac{\|\mathbf z\|^2}{R^2}.
		\]
		
		Since
		
		\[
		R=\|\mathbf z\|,
		\]
		
		it follows that
		
		\[
		\|\mathbf U\|^2
		=
		\frac{R^2}{R^2}
		=
		1.
		\]
		
		Therefore,
		
		\[
		\boxed{
			\|\mathbf U\|=1.
		}
		\]
		
		Hence $\mathbf U$ lies on the unit sphere
		
		\[
		S^\nu
		=
		\left\{
		x\in\mathbb R^{\nu+1}
		:
		\|x\|=1
		\right\}.
		\]
		
		\vspace{0.3cm}
		
		\paragraph{Step 3. Recover the Original Vector}
		
		Since
		
		\[
		\mathbf U
		=
		\frac{\mathbf z}{R},
		\]
		
		multiplying both sides by $R$ gives
		
		\[
		R\mathbf U
		=
		R
		\left(
		\frac{\mathbf z}{R}
		\right)
		=
		\mathbf z.
		\]
		
		Therefore every nonzero vector admits the decomposition
		
		\[
		\boxed{
			\mathbf z
			=
			R\mathbf U.
		}
		\]
		
		Geometrically,
		
		\begin{itemize}
			\item $R$ tells us \emph{how far} the point is from the origin.
			\item $\mathbf U$ tells us \emph{which direction} the point points toward.
		\end{itemize}
		
		Thus every vector is completely described by its radius and its direction.
		
		\vspace{0.3cm}
		
		\paragraph{Step 4. Passing to Random Vectors}
		
		The same construction applies to the random vector $\mathbf Z$.
		
		Define the random variables
		
		\[
		R=\|\mathbf Z\|,
		\]
		
		and
		
		\[
		\mathbf U
		=
		\frac{\mathbf Z}{\|\mathbf Z\|}.
		\]
		
		Then, with probability one,
		
		\[
		\boxed{
			\mathbf Z
			=
			R\mathbf U.
		}
		\]
		
		This is simply the polar-coordinate representation of a random vector.
		
		\vspace{0.3cm}
		
		\paragraph{Step 5. Why is $\mathbf U$ Uniform?}
		
		Recall that the density of $\mathbf Z$ depends only on the radius,
		
		\[
		f_{\mathbf Z}(\mathbf z)
		=
		\phi(\|\mathbf z\|).
		\]
		
		Consequently,
		
		\[
		f_{\mathbf Z}(Q\mathbf z)
		=
		f_{\mathbf Z}(\mathbf z)
		\]
		
		for every orthogonal matrix $Q$.
		
		In other words, rotating the vector does not change its probability
		distribution.
		
		Since rotations preserve the radius,
		
		\[
		\|Q\mathbf z\|
		=
		\|\mathbf z\|,
		\]
		
		the probability assigned to a point depends only on its distance from the
		origin and not on its orientation.
		
		Therefore, once the radius $R=r$ has been fixed, every point on the sphere
		
		\[
		\{\mathbf z:\|\mathbf z\|=r\}
		\]
		
		is equally likely.
		
		After dividing by the radius,
		
		\[
		\mathbf U=\frac{\mathbf Z}{R},
		\]
		
		every point on the unit sphere is equally likely.
		
		Thus,
		
		\[
		\boxed{
			\mathbf U
			\sim
			\text{Uniform}(S^\nu).
		}
		\]
		
		\vspace{0.3cm}
		
		\paragraph{Step 6. Independence of Radius and Direction}
		
		The rotational invariance of $\mathbf Z$ implies that changing the direction
		of the vector does not affect the distribution of its length.
		
		Likewise, changing the radius does not alter the distribution of the
		direction.
		
		Consequently, the random variables
		
		\[
		R
		\quad\text{and}\quad
		\mathbf U
		\]
		
		are independent.
		
		Therefore,
		
		\[
		\boxed{
			\mathbf Z
			=
			R\mathbf U,
			\qquad
			R=\|\mathbf Z\|,
			\qquad
			\mathbf U\sim\text{Uniform}(S^\nu),
			\qquad
			R\perp\!\!\!\perp\mathbf U.
		}
		\]
		
		This result is known as the \emph{polar representation theorem} for
		rotationally invariant random vectors. It separates the randomness into two
		independent components:
		
		\begin{enumerate}
			\item the random distance from the origin ($R$), and
			\item the random direction on the sphere ($\mathbf U$).
		\end{enumerate}
		
		This decomposition is the key step in the geometric derivation of the
		Student's $t$ distribution.
		
	\end{mdframed}
	
	Moreover,
	
	\[
	R
	\perp\!\!\!\perp
	\mathbf U.
	\]
	
	Thus,
	
	\[
	Z_i
	=
	RU_i,
	\qquad
	i=1,\ldots,\nu+1.
	\]
	
	%%%%%%%%%%%%%%%%%%%%%%%%%%%%%%%%%%%%%%%%%%%%%%%%%%%%%%%%%%%%%%%
	
	\section{Step 3. Rewriting the Classical Definition}
	
	\begin{mdframed}[
		linewidth=0.8pt,
		roundcorner=6pt,
		backgroundcolor=gray!5,
		linecolor=gray!60,
		innertopmargin=10pt,
		innerbottommargin=10pt,
		innerleftmargin=10pt,
		innerrightmargin=10pt
		]
		
		\subsection*{Why Does the Sample Variance Produce a Chi-Square Distribution?}
		
		The sample variance is defined by
		
		\[
		S^2
		=
		\frac1{n-1}
		\sum_{i=1}^{n}
		(X_i-\bar X)^2.
		\]
		
		Unlike the sample mean, its distribution is not immediately obvious.
		
		To understand it, first rewrite the definition as
		
		\[
		(n-1)S^2
		=
		\sum_{i=1}^{n}
		(X_i-\bar X)^2.
		\]
		
		Dividing both sides by $\sigma^2$ gives
		
		\[
		\frac{(n-1)S^2}{\sigma^2}
		=
		\sum_{i=1}^{n}
		\left(
		\frac{X_i-\bar X}{\sigma}
		\right)^2.
		\]
		
		At first glance, one might expect this to be the sum of $n$
		independent squared standard normal random variables.
		
		However, this is not true.
		
		The quantities
		
		\[
		X_i-\bar X
		\]
		
		are not independent because they satisfy
		
		\[
		\sum_{i=1}^{n}
		(X_i-\bar X)=0.
		\]
		
		This equation imposes one linear constraint on the deviations from the
		sample mean.
		
		Consequently, only
		
		\[
		n-1
		\]
		
		of these quantities can vary independently.
		
		The remarkable fact is given by the following theorem.
		
		\paragraph{Normal Sampling Theorem (Cochran's Theorem)}
		
		If
		
		\[
		X_1,\ldots,X_n
		\]
		
		are independent observations from
		
		\[
		N(\mu,\sigma^2),
		\]
		
		then
		
		\[
		\boxed{
			\frac{(n-1)S^2}{\sigma^2}
			\sim
			\chi^2_{\,n-1}.
		}
		\]
		
		Moreover,
		
		\[
		\boxed{
			\bar X
			\quad\text{and}\quad
			S^2
			\text{ are independent.}
		}
		\]
		
		The reduction from $n$ to $n-1$ degrees of freedom occurs because one degree
		of freedom is used to estimate the unknown population mean.
		
		Once the sample mean has been computed, only $n-1$ independent deviations
		remain.
		
		This result is highly special to the normal distribution and generally
		fails for other probability distributions.
		
		It is precisely this theorem that makes the Student's $t$ distribution
		possible, since it allows the statistic
		
		\[
		\frac{\bar X-\mu}{S/\sqrt n}
		\]
		
		to be written as the ratio of
		
		\begin{itemize}
			\item a standard normal random variable, and
			\item the square root of an independent chi-square random variable divided
			by its degrees of freedom.
		\end{itemize}
		
	\end{mdframed}
	
	\begin{mdframed}[
		linewidth=0.8pt,
		roundcorner=6pt,
		backgroundcolor=gray!5,
		linecolor=gray!60,
		innertopmargin=10pt,
		innerbottommargin=10pt,
		innerleftmargin=10pt,
		innerrightmargin=10pt
		]
		
		\subsection*{Why Does the Student's $t$ Statistic Arise?}
		
		The Student's $t$ statistic was introduced by William Sealy Gosset in 1908
		while working at the Guinness Brewery. The motivating problem was not the
		study of geometric objects or probability distributions, but rather a
		fundamental question in statistical inference:
		
		\begin{center}
			\emph{How can we estimate the mean of a normal population when the
				population standard deviation is unknown?}
		\end{center}
		
		Suppose
		
		\[
		X_1,X_2,\ldots,X_n
		\]
		
		are independent observations from a normal distribution
		
		\[
		X_i\sim N(\mu,\sigma^2),
		\]
		
		where both the population mean $\mu$ and the population standard deviation
		$\sigma$ are unknown.
		
		Our goal is to estimate $\mu$ and construct confidence intervals or perform
		hypothesis tests about its value.
		
		\paragraph{Step 1. The Sample Mean}
		
		The natural estimator of the population mean is the sample mean
		
		\[
		\bar X
		=
		\frac1n
		\sum_{i=1}^{n}X_i.
		\]
		
		Since the observations are independent and normally distributed,
		
		\[
		\bar X
		\sim
		N\left(
		\mu,
		\frac{\sigma^2}{n}
		\right).
		\]
		
		Thus,
		
		\[
		E[\bar X]=\mu,
		\]
		
		and
		
		\[
		\operatorname{Var}(\bar X)
		=
		\frac{\sigma^2}{n}.
		\]
		
		The standard deviation of the sample mean is therefore
		
		\[
		\frac{\sigma}{\sqrt n},
		\]
		
		which is known as the \emph{standard error} of the sample mean.
		
		\paragraph{Step 2. Standardizing the Sample Mean}
		
		Suppose, temporarily, that the population standard deviation $\sigma$ were
		known.
		
		Then we could standardize the sample mean by subtracting its expectation
		and dividing by its standard deviation:
		
		\[
		Z
		=
		\frac{\bar X-\mu}
		{\sigma/\sqrt n}.
		\]
		
		Since $\bar X$ is normally distributed,
		
		\[
		Z\sim N(0,1).
		\]
		
		This quantity is called the \emph{$z$-statistic} or \emph{$z$-score}.
		
		Because its distribution is completely known, probabilities involving $Z$
		can be computed exactly. Consequently, confidence intervals and hypothesis
		tests for $\mu$ are straightforward to construct.
		
		\paragraph{Step 3. The Practical Problem}
		
		Unfortunately, in almost every real application the population standard
		deviation $\sigma$ is unknown.
		
		Therefore, the statistic
		
		\[
		\frac{\bar X-\mu}
		{\sigma/\sqrt n}
		\]
		
		cannot actually be computed.
		
		Instead, we estimate $\sigma$ using the sample standard deviation,
		
		\[
		S
		=
		\sqrt{
			\frac1{n-1}
			\sum_{i=1}^{n}
			(X_i-\bar X)^2
		}.
		\]
		
		Replacing $\sigma$ by $S$ gives
		
		\[
		\frac{\bar X-\mu}
		{S/\sqrt n}.
		\]
		
		This seems like a natural substitution, but it creates a new problem.
		
		The denominator is now random.
		
		Consequently, the resulting statistic is no longer normally distributed.
		
		\paragraph{Step 4. Determining the New Distribution}
		
		To determine the distribution of
		
		\[
		\frac{\bar X-\mu}
		{S/\sqrt n},
		\]
		
		we use two remarkable properties that hold only for samples drawn from a
		normal distribution.
		
		First,
		
		\[
		\frac{\bar X-\mu}
		{\sigma/\sqrt n}
		\sim
		N(0,1).
		\]
		
		Second,
		
		\[
		\frac{(n-1)S^2}{\sigma^2}
		\sim
		\chi^2_{n-1}.
		\]
		
		Even more importantly,
		
		\[
		\bar X
		\quad\text{and}\quad
		S^2
		\]
		
		are independent random variables.
		
		This independence is a special property of the normal distribution and does
		not generally hold for other distributions.
		
		Therefore,
		
		\[
		\frac{\bar X-\mu}
		{S/\sqrt n}
		=
		\frac{
			\dfrac{\bar X-\mu}{\sigma/\sqrt n}
		}
		{
			\sqrt{
				\dfrac{(n-1)S^2/\sigma^2}{n-1}
			}
		}.
		\]
		
		The numerator is a standard normal random variable, while the denominator
		is the square root of an independent chi-square random variable divided by
		its degrees of freedom.
		
		Thus, if we define
		
		\[
		Z
		=
		\frac{\bar X-\mu}
		{\sigma/\sqrt n},
		\]
		
		and
		
		\[
		W
		=
		\frac{(n-1)S^2}{\sigma^2},
		\]
		
		then
		
		\[
		Z\sim N(0,1),
		\qquad
		W\sim\chi^2_{n-1},
		\]
		
		with
		
		\[
		Z
		\perp\!\!\!\perp
		W.
		\]
		
		Hence,
		
		\[
		\boxed{
			T
			=
			\frac{Z}
			{\sqrt{W/(n-1)}}
			=
			\frac{\bar X-\mu}
			{S/\sqrt n}.
		}
		\]
		
		This random variable follows the Student's $t$ distribution with
		
		\[
		n-1
		\]
		
		degrees of freedom.
		
		\paragraph{Why This Definition?}
		
		The definition of the Student's $t$ statistic is therefore not arbitrary.
		
		It is the statistic that naturally appears when the unknown population
		standard deviation in the classical $z$-statistic is replaced by its sample
		estimate.
		
		The additional randomness introduced by estimating $\sigma$ changes the
		distribution from the standard normal distribution to the Student's
		$t$ distribution.
		
		As the sample size increases, the sample standard deviation becomes a more
		accurate estimate of the population standard deviation. Consequently,
		
		\[
		S
		\longrightarrow
		\sigma,
		\]
		
		and the Student's $t$ distribution approaches the standard normal
		distribution.
		
		Thus, the $t$ statistic can be viewed as the natural generalization of the
		classical $z$ statistic to the realistic setting in which the population
		variance is unknown.
		
	\end{mdframed}
	
	\begin{mdframed}[
		linewidth=0.8pt,
		roundcorner=6pt,
		backgroundcolor=gray!5,
		linecolor=gray!60,
		innertopmargin=10pt,
		innerbottommargin=10pt,
		innerleftmargin=10pt,
		innerrightmargin=10pt
		]
		
		\subsection*{Rewriting the Classical Definition of the Student's $t$ Distribution}
		
		The classical definition of the Student's $t$ statistic is
		
		\[
		T
		=
		\frac{Z_1}
		{\sqrt{
				\frac{
					Z_2^2+Z_3^2+\cdots+Z_{\nu+1}^2
				}{\nu}
		}},
		\]
		
		where
		
		\[
		Z_1,Z_2,\ldots,Z_{\nu+1}
		\]
		
		are independent standard normal random variables.
		
		Observe that the numerator consists only of the first coordinate of the
		Gaussian vector,
		
		\[
		\mathbf Z
		=
		(Z_1,\ldots,Z_{\nu+1})^T,
		\]
		
		whereas the denominator is the Euclidean length of the remaining
		$\nu$ coordinates divided by $\sqrt{\nu}$.
		
		Our goal is to express this quantity using the polar representation
		
		\[
		\mathbf Z
		=
		R\mathbf U,
		\]
		
		where
		
		\[
		R=\|\mathbf Z\|,
		\]
		
		and
		
		\[
		\mathbf U
		=
		(U_1,U_2,\ldots,U_{\nu+1})
		\]
		
		is uniformly distributed on the unit sphere.
		
		\vspace{0.3cm}
		
		\paragraph{Step 1. Substitute the Polar Representation}
		
		Since
		
		\[
		\mathbf Z=R\mathbf U,
		\]
		
		each coordinate satisfies
		
		\[
		Z_i=RU_i,
		\qquad
		i=1,\ldots,\nu+1.
		\]
		
		Substituting these expressions into the definition of $T$ gives
		
		\[
		T
		=
		\frac{RU_1}
		{
			\sqrt{
				\frac{
					(RU_2)^2+(RU_3)^2+\cdots+(RU_{\nu+1})^2
				}{\nu}
			}
		}.
		\]
		
		\vspace{0.3cm}
		
		\paragraph{Step 2. Expand the Squares}
		
		Since
		
		\[
		(RU_i)^2=R^2U_i^2,
		\]
		
		the denominator becomes
		
		\[
		\sqrt{
			\frac{
				R^2U_2^2
				+
				R^2U_3^2
				+\cdots+
				R^2U_{\nu+1}^2
			}{\nu}
		}.
		\]
		
		Factor the common term $R^2$:
		
		\[
		=
		\sqrt{
			\frac{
				R^2
				\left(
				U_2^2+\cdots+U_{\nu+1}^2
				\right)
			}{\nu}
		}.
		\]
		
		Thus,
		
		\[
		T
		=
		\frac{RU_1}
		{
			\sqrt{
				\frac{
					R^2
					\left(
					U_2^2+\cdots+U_{\nu+1}^2
					\right)
				}{\nu}
			}
		}.
		\]
		
		\vspace{0.3cm}
		
		\paragraph{Step 3. Simplify the Square Root}
		
		Use the elementary property
		
		\[
		\sqrt{ab}
		=
		\sqrt a\,\sqrt b.
		\]
		
		Since
		
		\[
		R\ge0,
		\]
		
		we have
		
		\[
		\sqrt{R^2}=R.
		\]
		
		Therefore,
		
		\[
		\sqrt{
			\frac{
				R^2
				\left(
				U_2^2+\cdots+U_{\nu+1}^2
				\right)
			}{\nu}
		}
		=
		R
		\sqrt{
			\frac{
				U_2^2+\cdots+U_{\nu+1}^2
			}{\nu}
		}.
		\]
		
		The expression for $T$ now becomes
		
		\[
		T
		=
		\frac{RU_1}
		{
			R
			\sqrt{
				\frac{
					U_2^2+\cdots+U_{\nu+1}^2
				}{\nu}
			}
		}.
		\]
		
		\vspace{0.3cm}
		
		\paragraph{Step 4. Cancel the Radius}
		
		Since the factor $R$ appears in both the numerator and denominator,
		
		\[
		T
		=
		\frac{U_1}
		{
			\sqrt{
				\frac{
					U_2^2+\cdots+U_{\nu+1}^2
				}{\nu}
			}
		}.
		\]
		
		This is the crucial observation of the entire derivation.
		
		The random radius has completely disappeared.
		
		Therefore, the Student's $t$ statistic depends only on the direction of the
		Gaussian vector and not on its distance from the origin.
		
		\vspace{0.3cm}
		
		\paragraph{Step 5. Use the Unit Sphere Constraint}
		
		Because
		
		\[
		\mathbf U
		\]
		
		lies on the unit sphere,
		
		\[
		\|\mathbf U\|=1.
		\]
		
		Hence,
		
		\[
		U_1^2+U_2^2+\cdots+U_{\nu+1}^2=1.
		\]
		
		Solving for the remaining coordinates,
		
		\[
		U_2^2+\cdots+U_{\nu+1}^2
		=
		1-U_1^2.
		\]
		
		Substituting this identity into the denominator,
		
		\[
		T
		=
		\frac{U_1}
		{
			\sqrt{
				\frac{
					1-U_1^2
				}{\nu}
			}
		}.
		\]
		
		\vspace{0.3cm}
		
		\paragraph{Step 6. Move the Factor $\sqrt{\nu}$ Outside}
		
		Using
		
		\[
		\sqrt{\frac{a}{\nu}}
		=
		\frac{\sqrt a}{\sqrt\nu},
		\]
		
		we obtain
		
		\[
		T
		=
		\frac{U_1}
		{
			\dfrac{\sqrt{1-U_1^2}}{\sqrt\nu}
		}.
		\]
		
		Dividing by a fraction is equivalent to multiplying by its reciprocal.
		
		Therefore,
		
		\[
		T
		=
		U_1
		\cdot
		\frac{\sqrt\nu}
		{\sqrt{1-U_1^2}}.
		\]
		
		Finally,
		
		\[
		\boxed{
			T
			=
			\sqrt\nu\,
			\frac{U_1}
			{\sqrt{1-U_1^2}}.
		}
		\]
		
		This formula reveals the geometric nature of the Student's $t$
		distribution. Every dependence on the Gaussian radius has vanished; the
		distribution depends only on the first coordinate of a uniformly random
		point on the unit sphere. In the next step, we interpret this coordinate
		as the cosine of the angle between the random direction $\mathbf U$ and a
		fixed reference axis.
		
	\end{mdframed}
	
	Substitute
	
	\[
	Z_i=RU_i.
	\]
	
	Then
	
	\[
	T
	=
	\frac{RU_1}
	{
		\sqrt{
			\frac{
				R^2(U_2^2+\cdots+U_{\nu+1}^2)
			}{\nu}
		}
	}.
	\]
	
	Since
	
	\[
	\sqrt{R^2}=R,
	\]
	
	the radius cancels:
	
	\[
	T
	=
	\frac{U_1}
	{
		\sqrt{
			\frac{
				U_2^2+\cdots+U_{\nu+1}^2
			}{\nu}
		}
	}.
	\]
	
	Now use the defining property of the unit sphere,
	
	\[
	U_1^2+\cdots+U_{\nu+1}^2=1.
	\]
	
	Therefore
	
	\[
	U_2^2+\cdots+U_{\nu+1}^2
	=
	1-U_1^2.
	\]
	
	Substituting,
	
	\[
	T
	=
	\frac{U_1}
	{
		\sqrt{
			\frac{1-U_1^2}{\nu}
		}
	}.
	\]
	
	Finally,
	
	\[
	\boxed{
		T
		=
		\sqrt{\nu}
		\frac{U_1}{\sqrt{1-U_1^2}}
	}
	\]
	
	which no longer contains any Gaussian random variables.
	
	%%%%%%%%%%%%%%%%%%%%%%%%%%%%%%%%%%%%%%%%%%%%%%%%%%%%%%%%%%%%%%%
	
	\section{Step 4. Introducing the Angle}
	
	Fix any unit vector
	
	\[
	\mathbf e_1=(1,0,\ldots,0).
	\]
	
	Let
	
	\[
	\theta
	\]
	
	be the angle between $\mathbf U$ and $\mathbf e_1$.
	
	By the definition of the dot product,
	
	\[
	\cos\theta
	=
	\mathbf U\cdot\mathbf e_1
	=
	U_1.
	\]
	
	Hence,
	
	\[
	1-U_1^2
	=
	1-\cos^2\theta
	=
	\sin^2\theta.
	\]
	
	Therefore,
	
	\[
	\frac{U_1}{\sqrt{1-U_1^2}}
	=
	\frac{\cos\theta}{\sin\theta}
	=
	\cot\theta.
	\]
	
	Thus,
	
	\[
	\boxed{
		T
		=
		\sqrt{\nu}\cot\theta.
	}
	\]
	
	This is the geometric interpretation of the Student's
	$t$ distribution.
	
	%%%%%%%%%%%%%%%%%%%%%%%%%%%%%%%%%%%%%%%%%%%%%%%%%%%%%%%%%%%%%%%
	
	\section{Step 5. Finding the Distribution of the Angle}
	
	Suppose
	
	\[
	\theta
	\]
	
	is fixed.
	
	Then
	
	\[
	U_1=\cos\theta.
	\]
	
	Since
	
	\[
	U_1^2+\cdots+U_{\nu+1}^2=1,
	\]
	
	we obtain
	
	\[
	U_2^2+\cdots+U_{\nu+1}^2
	=
	\sin^2\theta.
	\]
	
	This equation describes a sphere of
	
	\[
	(\nu-1)
	\]
	
	dimensions with radius
	
	\[
	\sin\theta.
	\]
	
	%%%%%%%%%%%%%%%%%%%%%%%%%%%%%%%%%%%%%%%%%%%%%%%%%%%%%%%%%%%%%%%
	
	\section{Step 6. Surface Area of the Latitude}
	
	The surface area of a $d$-dimensional sphere scales like
	
	\[
	r^d.
	\]
	
	Therefore,
	
	the sphere obtained by fixing $\theta$ has area proportional to
	
	\[
	(\sin\theta)^{\nu-1}.
	\]
	
	Now consider a thin strip between
	
	\[
	\theta
	\quad\text{and}\quad
	\theta+d\theta.
	\]
	
	Its area is approximately
	
	\[
	(\sin\theta)^{\nu-1}
	\,d\theta.
	\]
	
	Since points are uniformly distributed,
	
	\[
	P(\theta\in[\theta,\theta+d\theta])
	\propto
	(\sin\theta)^{\nu-1}
	d\theta.
	\]
	
	Hence
	
	\[
	f_\Theta(\theta)
	=
	C
	(\sin\theta)^{\nu-1},
	\]
	
	where $C$ is the normalizing constant.
	
	%%%%%%%%%%%%%%%%%%%%%%%%%%%%%%%%%%%%%%%%%%%%%%%%%%%%%%%%%%%%%%%
	
	\section{Step 7. Computing the Constant}
	
	The density must integrate to one:
	
	\[
	1
	=
	C
	\int_0^\pi
	(\sin\theta)^{\nu-1}
	d\theta.
	\]
	
	Using
	
	\[
	\int_0^\pi
	(\sin\theta)^{\nu-1}
	d\theta
	=
	B
	\left(
	\frac12,
	\frac{\nu}{2}
	\right),
	\]
	
	we obtain
	
	\[
	C
	=
	\frac1{
		B\left(\frac12,\frac\nu2\right)
	}.
	\]
	
	Therefore,
	
	\[
	\boxed{
		f_\Theta(\theta)
		=
		\frac{
			(\sin\theta)^{\nu-1}
		}{
			B\left(\frac12,\frac\nu2\right)
		},
		\qquad
		0<\theta<\pi.
	}
	\]
	
	%%%%%%%%%%%%%%%%%%%%%%%%%%%%%%%%%%%%%%%%%%%%%%%%%%%%%%%%%%%%%%%
	
	\section{Step 8. Transforming to the $t$ Distribution}
	
	Recall
	
	\[
	T
	=
	\sqrt{\nu}\cot\theta.
	\]
	
	Differentiate:
	
	\[
	\frac{dT}{d\theta}
	=
	-\sqrt{\nu}\csc^2\theta.
	\]
	
	Therefore,
	
	\[
	\left|
	\frac{d\theta}{dT}
	\right|
	=
	\frac{\sin^2\theta}{\sqrt{\nu}}.
	\]
	
	The change-of-variable formula gives
	
	\[
	f_T(t)
	=
	f_\Theta(\theta)
	\left|
	\frac{d\theta}{dt}
	\right|.
	\]
	
	Substitute the expressions obtained above:
	
	\[
	f_T(t)
	=
	\frac{
		(\sin\theta)^{\nu-1}
	}{
		B\left(\frac12,\frac\nu2\right)
	}
	\cdot
	\frac{\sin^2\theta}{\sqrt{\nu}}.
	\]
	
	Therefore,
	
	\[
	f_T(t)
	=
	\frac{
		(\sin\theta)^{\nu+1}
	}{
		\sqrt{\nu}
		B\left(\frac12,\frac\nu2\right)
	}.
	\]
	
	%%%%%%%%%%%%%%%%%%%%%%%%%%%%%%%%%%%%%%%%%%%%%%%%%%%%%%%%%%%%%%%
	
	\section{Step 9. Expressing Everything in Terms of $t$}
	
	Since
	
	\[
	t
	=
	\sqrt{\nu}\cot\theta,
	\]
	
	we have
	
	\[
	\cot\theta
	=
	\frac{t}{\sqrt{\nu}}.
	\]
	
	Using
	
	\[
	1+\cot^2\theta
	=
	\csc^2\theta,
	\]
	
	we obtain
	
	\[
	\sin^2\theta
	=
	\frac1{1+\cot^2\theta}
	=
	\frac1{1+t^2/\nu}
	=
	\frac{\nu}{\nu+t^2}.
	\]
	
	Hence,
	
	\[
	\sin\theta
	=
	\sqrt{\frac{\nu}{\nu+t^2}}.
	\]
	
	Substitute into the density:
	
	\[
	f_T(t)
	=
	\frac1{
		\sqrt{\nu}
		B\left(\frac12,\frac\nu2\right)
	}
	\left(
	\frac{\nu}{\nu+t^2}
	\right)^{(\nu+1)/2}.
	\]
	
	Rewrite as
	
	\[
	\boxed{
		f_T(t)
		=
		\frac1{
			\sqrt{\nu}
			B\left(\frac12,\frac\nu2\right)
		}
		\left(
		1+\frac{t^2}{\nu}
		\right)^{-(\nu+1)/2}.
	}
	\]
	
	Finally, using
	
	\[
	B(a,b)
	=
	\frac{\Gamma(a)\Gamma(b)}
	{\Gamma(a+b)}
	\]
	
	and
	
	\[
	\Gamma\!\left(\frac12\right)
	=
	\sqrt{\pi},
	\]
	
	we recover the familiar form
	
	\[
	\boxed{
		f_T(t)
		=
		\frac{
			\Gamma\left(\frac{\nu+1}{2}\right)
		}{
			\sqrt{\nu\pi}\,
			\Gamma\left(\frac\nu2\right)
		}
		\left(
		1+\frac{t^2}{\nu}
		\right)^{-(\nu+1)/2}.
	}
	\]
	
	%%%%%%%%%%%%%%%%%%%%%%%%%%%%%%%%%%%%%%%%%%%%%%%%%%%%%%%%%%%%%%%
	
	\section{Conclusion}
	
	The Student's $t$ distribution is fundamentally a geometric object.
	
	Its origin is not the Gaussian distribution itself, but the geometry of uniformly distributed points on high-dimensional spheres. The Gaussian derivation merely provides one convenient construction because Gaussian vectors are rotationally invariant.
	
	The entire derivation can be summarized as
	
	\[
	\boxed{
		\text{Uniform point on }S^\nu
		\Longrightarrow
		\theta
		\Longrightarrow
		T=\sqrt{\nu}\cot\theta
		\Longrightarrow
		\text{Student's }t.
	}
	\]
	
\end{document}